\section{个人路径：从 App Store 到 Agent 创业}

\subsection{第一阶段：移动互联网红利期的产品直觉}

季逸超高中时期做第三方 iOS 浏览器并实现几十万美金收入，这段经历塑造了两个后续长期有效的认知：第一，用户愿意为“速度和省心”付费；第二，技术窗口期稍纵即逝，速度本身就是竞争力。早期的 Buy Copy 模式虽然简单，但训练了最难得的能力：把技术决策直接映射到现金流。

\begin{importantbox}{早期创业训练出的硬能力}
\begin{itemize}
    \item 把“用户抱怨”翻译成工程需求，而不是写成抽象愿景
    \item 在算力和人力都有限时，优先解决“最高频痛点”
    \item 用简单可验证的指标替代复杂 KPI
\end{itemize}
\end{importantbox}

\subsection{第二阶段：从预加载需求转向 NLP 研究}

从浏览器业务过渡到 NLP 的触发因素并非“学术兴趣”，而是产品问题：如何预测下一步点击并提前加载。这个问题本质上是序列建模问题。Word2Vec 发布后，文本稠密表示可用，工程团队终于有机会把经验规则升级为可泛化的模型能力。

\begin{knowledgebox}{技术迁移的实际代价}
访谈里反复提到一件事：技术代际切换不会给旧系统“平滑升级”机会。Word2Vec 到 LSTM，再到 Transformer/BERT，每一轮都需要重写一部分数据管线、标注规范与训练脚本。
\end{knowledgebox}

\subsection{第三阶段：知识图谱创业与失败复盘}

第二次创业聚焦语义搜索和 Open IE，目标是突破传统“十个蓝色链接”范式。团队做了模型预训练、向量索引、图谱构建等大量底层工作，但产品层面没有形成持续增长。访谈中的复盘非常直接：技术领先不等于商业成立，尤其在生态、渠道和数据授权都不占优时。

\begin{warningbox}{“太早”的三种具体表现}
\begin{itemize}
    \item 交互入口尚未成熟：可穿戴设备和语音界面未形成主流分发
    \item 非技术壁垒被低估：数据供给与合作关系决定了上限
    \item 商业模式摇摆：2C 与 2B 切换导致团队组织目标撕裂
\end{itemize}
\end{warningbox}

\subsection{第四阶段：企业内 LLM 实战与再创业决心}

加入独角兽公司后，季逸超在“打榜体系”中系统化了 Benchmark 思维：如何把需求变成可比较指标，如何让资源分配与评测结果对齐。GPT-3 之后，行业范式切换，他再次看到创业窗口。与第一次创业相比，这次更强调“系统工程 + 产品闭环”，而不是单点模型能力。

\begin{importantbox}{从“做模型”到“做系统”}
访谈中最值得记的变化是目标函数变化：过去关注模型指标，后来关注任务闭环完成率（task completion reliability）。这决定了 Manus 选择了“模型 + 沙盒 + 工具”的系统路线。
\end{importantbox}

\subsection{本章小结}

季逸超的路径不是线性升级，而是多次“问题驱动转向”。真正可迁移的不是某个模型栈，而是三个动作：快速验证、及时止损、把失败沉淀为下一阶段的边界条件。


